% Hand-written bibliography, so that no BibTeX run is needed. Every journal
% entry was checked against the publisher-deposited Crossref record at the DOI
% given; the six Protein Data Bank primary citations were read from the RCSB
% entry records and then checked at Crossref. Software without a paper is
% cited as software with a URL and an access date. Order follows first
% appearance in the text.
\begin{thebibliography}{99}
\small
\setlength{\itemsep}{1pt}

\bibitem{dror2012microscope}
R.~O. Dror, R.~M. Dirks, J.~P. Grossman, H.~Xu, and D.~E. Shaw,
``Biomolecular simulation: a computational microscope for molecular biology,''
\emph{Annu. Rev. Biophys.} \textbf{41}(1), 429--452 (2012).
DOI: 10.1146/annurev-biophys-042910-155245.

\bibitem{shaw2010anton}
D.~E. Shaw, P.~Maragakis, K.~Lindorff-Larsen, S.~Piana, R.~O. Dror,
M.~P. Eastwood, J.~A. Bank, J.~M. Jumper, J.~K. Salmon, Y.~Shan, and
W.~Wriggers,
``Atomic-level characterization of the structural dynamics of proteins,''
\emph{Science} \textbf{330}(6002), 341--346 (2010).
DOI: 10.1126/science.1187409.

\bibitem{lindorfflarsen2011folding}
K.~Lindorff-Larsen, S.~Piana, R.~O. Dror, and D.~E. Shaw,
``How fast-folding proteins fold,''
\emph{Science} \textbf{334}(6055), 517--520 (2011).
DOI: 10.1126/science.1208351.

\bibitem{torrie1977umbrella}
G.~M. Torrie and J.~P. Valleau,
``Nonphysical sampling distributions in Monte Carlo free-energy estimation:
umbrella sampling,''
\emph{J. Comput. Phys.} \textbf{23}(2), 187--199 (1977).
DOI: 10.1016/0021-9991(77)90121-8.

\bibitem{sugita1999remd}
Y.~Sugita and Y.~Okamoto,
``Replica-exchange molecular dynamics method for protein folding,''
\emph{Chem. Phys. Lett.} \textbf{314}(1--2), 141--151 (1999).
DOI: 10.1016/S0009-2614(99)01123-9.

\bibitem{laio2002metadynamics}
A.~Laio and M.~Parrinello,
``Escaping free-energy minima,''
\emph{Proc. Natl. Acad. Sci. U.S.A.} \textbf{99}(20), 12562--12566 (2002).
DOI: 10.1073/pnas.202427399.

\bibitem{bernardi2015enhanced}
R.~C. Bernardi, M.~C.~R. Melo, and K.~Schulten,
``Enhanced sampling techniques in molecular dynamics simulations of biological
systems,''
\emph{Biochim. Biophys. Acta Gen. Subj.} \textbf{1850}(5), 872--877 (2015).
DOI: 10.1016/j.bbagen.2014.10.019.

\bibitem{minary2010chainclosure}
P.~Minary and M.~Levitt,
``Conformational optimization with natural degrees of freedom: a novel
stochastic chain closure algorithm,''
\emph{J. Comput. Biol.} \textbf{17}(8), 993--1010 (2010).
DOI: 10.1089/cmb.2010.0016.

\bibitem{sim2012naturalmoves}
A.~Y.~L. Sim, M.~Levitt, and P.~Minary,
``Modeling and design by hierarchical natural moves,''
\emph{Proc. Natl. Acad. Sci. U.S.A.} \textbf{109}(8), 2890--2895 (2012).
DOI: 10.1073/pnas.1119918109.

\bibitem{minary2014nucleosome}
P.~Minary and M.~Levitt,
``Training-free atomistic prediction of nucleosome occupancy,''
\emph{Proc. Natl. Acad. Sci. U.S.A.} \textbf{111}(17), 6293--6298 (2014).
DOI: 10.1073/pnas.1404475111.

\bibitem{zhang2012multiscale}
J.~Zhang, P.~Minary, and M.~Levitt,
``Multiscale natural moves refine macromolecules using single-particle electron
microscopy projection images,''
\emph{Proc. Natl. Acad. Sci. U.S.A.} \textbf{109}(25), 9845--9850 (2012).
DOI: 10.1073/pnas.1205945109.

\bibitem{chang2020rnacryoem}
J.-Y. Chang, Z.~Cui, K.~Yang, J.~Huang, P.~Minary, and J.~Zhang,
``Hierarchical natural move Monte Carlo refines flexible RNA structures into
cryo-EM densities,''
\emph{RNA} \textbf{26}(12), 1755--1766 (2020).
DOI: 10.1261/rna.071100.119.

\bibitem{krawczyk2018epigenetic}
K.~Krawczyk, S.~Demharter, B.~Knapp, C.~M. Deane, and P.~Minary,
``In silico structural modeling of multiple epigenetic marks on DNA,''
\emph{Bioinformatics} \textbf{34}(1), 41--48 (2018).
DOI: 10.1093/bioinformatics/btx516.

\bibitem{knapp2016pmhc}
B.~Knapp, S.~Demharter, C.~M. Deane, and P.~Minary,
``Exploring peptide/MHC detachment processes using hierarchical natural move
Monte Carlo,''
\emph{Bioinformatics} \textbf{32}(2), 181--186 (2016).
DOI: 10.1093/bioinformatics/btv502.

\bibitem{vitalis2009mcmethods}
A.~Vitalis and R.~V. Pappu,
``Methods for Monte Carlo simulations of biomacromolecules,''
\emph{Annu. Rep. Comput. Chem.} \textbf{5}, 49--76 (2009).
DOI: 10.1016/S1574-1400(09)00503-9.

\bibitem{jorgensen2005mcpro}
W.~L. Jorgensen and J.~Tirado-Rives,
``Molecular modeling of organic and biomolecular systems using BOSS and
MCPRO,''
\emph{J. Comput. Chem.} \textbf{26}(16), 1689--1700 (2005).
DOI: 10.1002/jcc.20297.

\bibitem{protoms-software}
\emph{ProtoMS}, version 3.4.0, C.~J. Woods, J.~Michel, M.~Bodnarchuk,
S.~Genheden, R.~Bradshaw, G.~A. Ross, C.~Cave-Ayland, A.~I. Cabedo~Martinez,
H.~Bruce-Macdonald, J.~Graham, and M.~Samways, School of Chemistry, University
of Southampton, released under the GNU GPL (2018).
\url{https://www.essexgroup.soton.ac.uk/ProtoMS/} (accessed 25 August 2026).
The distribution names no method paper of its own and carries no DOI, so it is
cited as software.

\bibitem{shah2017cassandra}
J.~K. Shah, E.~Mar\'in-Rimoldi, R.~G. Mullen, B.~P. Keene, S.~Khan,
A.~S. Paluch, N.~Rai, L.~L. Romanielo, T.~W. Rosch, B.~Yoo, and E.~J. Maginn,
``Cassandra: an open source Monte Carlo package for molecular simulation,''
\emph{J. Comput. Chem.} \textbf{38}(19), 1727--1739 (2017).
DOI: 10.1002/jcc.24807.

\bibitem{ulmschneider2003concerted}
J.~P. Ulmschneider and W.~L. Jorgensen,
``Monte Carlo backbone sampling for polypeptides with variable bond angles and
dihedral angles using concerted rotations and a Gaussian bias,''
\emph{J. Chem. Phys.} \textbf{118}(9), 4261--4271 (2003).
DOI: 10.1063/1.1542611.

\bibitem{ulmschneider2004concertedna}
J.~P. Ulmschneider and W.~L. Jorgensen,
``Monte Carlo backbone sampling for nucleic acids using concerted rotations
including variable bond angles,''
\emph{J. Phys. Chem. B} \textbf{108}(43), 16883--16892 (2004).
DOI: 10.1021/jp047796z.

\bibitem{minary2008foldspace}
P.~Minary and M.~Levitt,
``Probing protein fold space with a simplified model,''
\emph{J. Mol. Biol.} \textbf{375}(4), 920--933 (2008).
DOI: 10.1016/j.jmb.2007.10.087.

\bibitem{leaverfay2011rosetta}
A.~Leaver-Fay, M.~Tyka, S.~M. Lewis, O.~F. Lange, J.~Thompson, R.~Jacak,
K.~W. Kaufman, P.~D. Renfrew, C.~A. Smith, W.~Sheffler, I.~W. Davis, S.~Cooper,
A.~Treuille, D.~J. Mandell, F.~Richter, Y.-E.~A. Ban, S.~J. Fleishman,
J.~E. Corn, D.~E. Kim, S.~Lyskov, M.~Berrondo, S.~Mentzer, Z.~Popovi\'c,
J.~J. Havranek, J.~Karanicolas, R.~Das, J.~Meiler, T.~Kortemme, J.~J. Gray,
B.~Kuhlman, D.~Baker, and P.~Bradley,
``ROSETTA3: an object-oriented software suite for the simulation and design of
macromolecules,''
\emph{Methods Enzymol.} \textbf{487}, 545--574 (2011).
DOI: 10.1016/B978-0-12-381270-4.00019-6.

\bibitem{pymosaics-plugin}
K.~Krawczyk, \emph{PymoSAICS}, a PyMOL plugin for MOSAICS, Department of
Computer Science, University of Oxford (2015).
\url{https://www.cs.ox.ac.uk/mosaics/} (accessed 6 September 2026). No
peer-reviewed publication describes the plugin, so it is cited as software.

\bibitem{maier2015ff14sb}
J.~A. Maier, C.~Martinez, K.~Kasavajhala, L.~Wickstrom, K.~E. Hauser, and
C.~Simmerling,
``ff14SB: improving the accuracy of protein side chain and backbone parameters
from ff99SB,''
\emph{J. Chem. Theory Comput.} \textbf{11}(8), 3696--3713 (2015).
DOI: 10.1021/acs.jctc.5b00255.

\bibitem{tian2020ff19sb}
C.~Tian, K.~Kasavajhala, K.~A.~A. Belfon, L.~Raguette, H.~Huang, A.~N. Migues,
J.~Bickel, Y.~Wang, J.~Pincay, Q.~Wu, and C.~Simmerling,
``ff19SB: amino-acid-specific protein backbone parameters trained against
quantum mechanics energy surfaces in solution,''
\emph{J. Chem. Theory Comput.} \textbf{16}(1), 528--552 (2020).
DOI: 10.1021/acs.jctc.9b00591.

\bibitem{mackerell2004cmap}
A.~D. MacKerell,~Jr., M.~Feig, and C.~L. Brooks,~III,
``Improved treatment of the protein backbone in empirical force fields,''
\emph{J. Am. Chem. Soc.} \textbf{126}(3), 698--699 (2004).
DOI: 10.1021/ja036959e.

\bibitem{mackerell2004cmapfull}
A.~D. MacKerell,~Jr., M.~Feig, and C.~L. Brooks,~III,
``Extending the treatment of backbone energetics in protein force fields:
limitations of gas-phase quantum mechanics in reproducing protein
conformational distributions in molecular dynamics simulations,''
\emph{J. Comput. Chem.} \textbf{25}(11), 1400--1415 (2004).
DOI: 10.1002/jcc.20065.

\bibitem{perez2007bsc0}
A.~P\'erez, I.~March\'an, D.~Svozil, J.~\v{S}poner, T.~E. Cheatham,~III,
C.~A. Laughton, and M.~Orozco,
``Refinement of the AMBER force field for nucleic acids: improving the
description of $\alpha/\gamma$ conformers,''
\emph{Biophys. J.} \textbf{92}(11), 3817--3829 (2007).
DOI: 10.1529/biophysj.106.097782.

\bibitem{ivani2016bsc1}
I.~Ivani, P.~D. Dans, A.~Noy, A.~P\'erez, I.~Faustino, A.~Hospital, J.~Walther,
P.~Andrio, R.~Go\~ni, A.~Balaceanu, G.~Portella, F.~Battistini, J.~L. Gelp\'i,
C.~Gonz\'alez, M.~Vendruscolo, C.~A. Laughton, S.~A. Harris, D.~A. Case, and
M.~Orozco,
``Parmbsc1: a refined force field for DNA simulations,''
\emph{Nat. Methods} \textbf{13}(1), 55--58 (2016).
DOI: 10.1038/nmeth.3658.

\bibitem{zgarbova2011ol3}
M.~Zgarbov\'a, M.~Otyepka, J.~\v{S}poner, A.~Ml\'adek, P.~Bana\v{s},
T.~E. Cheatham,~III, and P.~Jure\v{c}ka,
``Refinement of the Cornell et al. nucleic acids force field based on reference
quantum chemical calculations of glycosidic torsion profiles,''
\emph{J. Chem. Theory Comput.} \textbf{7}(9), 2886--2902 (2011).
DOI: 10.1021/ct200162x.

\bibitem{zgarbova2015ol15}
M.~Zgarbov\'a, J.~\v{S}poner, M.~Otyepka, T.~E. Cheatham,~III,
R.~Galindo-Murillo, and P.~Jure\v{c}ka,
``Refinement of the sugar--phosphate backbone torsion beta for AMBER force
fields improves the description of Z- and B-DNA,''
\emph{J. Chem. Theory Comput.} \textbf{11}(12), 5723--5736 (2015).
DOI: 10.1021/acs.jctc.5b00716.

\bibitem{zgarbova2021ol21}
M.~Zgarbov\'a, J.~\v{S}poner, and P.~Jure\v{c}ka,
``Z-DNA as a touchstone for additive empirical force fields and a refinement of
the alpha/gamma DNA torsions for AMBER,''
\emph{J. Chem. Theory Comput.} \textbf{17}(10), 6292--6301 (2021).
DOI: 10.1021/acs.jctc.1c00697.

\bibitem{zgarbova2025ol24}
M.~Zgarbov\'a, J.~\v{S}poner, and P.~Jure\v{c}ka,
``Refinement of the sugar puckering torsion potential in the AMBER DNA force
field,''
\emph{J. Chem. Theory Comput.} \textbf{21}(2), 833--846 (2025).
DOI: 10.1021/acs.jctc.4c01100.

\bibitem{galindomurillo2016amberdna}
R.~Galindo-Murillo, J.~C. Robertson, M.~Zgarbov\'a, J.~\v{S}poner, M.~Otyepka,
P.~Jure\v{c}ka, and T.~E. Cheatham,~III,
``Assessing the current state of Amber force field modifications for DNA,''
\emph{J. Chem. Theory Comput.} \textbf{12}(8), 4114--4127 (2016).
DOI: 10.1021/acs.jctc.6b00186.

\bibitem{love2023amberdna}
O.~Love, R.~Galindo-Murillo, M.~Zgarbov\'a, J.~\v{S}poner, P.~Jure\v{c}ka, and
T.~E. Cheatham,~III,
``Assessing the current state of Amber force field modifications for
DNA---2023 edition,''
\emph{J. Chem. Theory Comput.} \textbf{19}(13), 4299--4307 (2023).
DOI: 10.1021/acs.jctc.3c00233.

\bibitem{hart2012charmm36dna}
K.~Hart, N.~Foloppe, C.~M. Baker, E.~J. Denning, L.~Nilsson, and
A.~D. MacKerell,~Jr.,
``Optimization of the CHARMM additive force field for DNA: improved treatment
of the BI/BII conformational equilibrium,''
\emph{J. Chem. Theory Comput.} \textbf{8}(1), 348--362 (2012).
DOI: 10.1021/ct200723y.

\bibitem{denning2011charmm36rna}
E.~J. Denning, U.~D. Priyakumar, L.~Nilsson, and A.~D. MacKerell,~Jr.,
``Impact of 2$'$-hydroxyl sampling on the conformational properties of RNA:
update of the CHARMM all-atom additive force field for RNA,''
\emph{J. Comput. Chem.} \textbf{32}(9), 1929--1943 (2011).
DOI: 10.1002/jcc.21777.

\bibitem{case2023ambertools}
D.~A. Case, H.~M. Aktulga, K.~Belfon, D.~S. Cerutti, G.~A. Cisneros,
V.~W.~D. Cruzeiro, N.~Forouzesh, T.~J. Giese, A.~W. G\"otz, H.~Gohlke,
S.~Izadi, K.~Kasavajhala, M.~C. Kaymak, E.~King, T.~Kurtzman, T.-S. Lee,
P.~Li, J.~Liu, T.~Luchko, R.~Luo, M.~Manathunga, M.~R. Machado, H.~M. Nguyen,
K.~A. O'Hearn, A.~V. Onufriev, F.~Pan, S.~Pantano, R.~Qi, A.~Rahnamoun,
A.~Risheh, S.~Schott-Verdugo, A.~Shajan, J.~Swails, J.~Wang, H.~Wei, X.~Wu,
Y.~Wu, S.~Zhang, S.~Zhao, Q.~Zhu, T.~E. Cheatham,~III, D.~R. Roe,
A.~Roitberg, C.~Simmerling, D.~M. York, M.~C. Nagan, and K.~M. Merz,~Jr.,
``AmberTools,''
\emph{J. Chem. Inf. Model.} \textbf{63}(20), 6183--6191 (2023).
DOI: 10.1021/acs.jcim.3c01153.
The calculations here used sander 24.0 from AmberTools 24.8.

\bibitem{eastman2017openmm7}
P.~Eastman, J.~Swails, J.~D. Chodera, R.~T. McGibbon, Y.~Zhao, K.~A. Beauchamp,
L.-P. Wang, A.~C. Simmonett, M.~P. Harrigan, C.~D. Stern, R.~P. Wiewiora,
B.~R. Brooks, and V.~S. Pande,
``OpenMM 7: rapid development of high performance algorithms for molecular
dynamics,''
\emph{PLoS Comput. Biol.} \textbf{13}(7), e1005659 (2017).
DOI: 10.1371/journal.pcbi.1005659.

\bibitem{eastman2024openmm8}
P.~Eastman, R.~Galvelis, R.~P. Pel\'aez, C.~R.~A. Abreu, S.~E. Farr,
E.~Gallicchio, A.~Gorenko, M.~M. Henry, F.~Hu, J.~Huang, A.~Kr\"amer,
J.~Michel, J.~A. Mitchell, V.~S. Pande, J.~P.~G.~L.~M. Rodrigues,
J.~Rodriguez-Guerra, A.~C. Simmonett, S.~Singh, J.~Swails, P.~Turner,
Y.~Wang, I.~Zhang, J.~D. Chodera, G.~De~Fabritiis, and T.~E. Markland,
``OpenMM 8: molecular dynamics simulation with machine learning potentials,''
\emph{J. Phys. Chem. B} \textbf{128}(1), 109--116 (2024).
DOI: 10.1021/acs.jpcb.3c06662.

\bibitem{abraham2015gromacs}
M.~J. Abraham, T.~Murtola, R.~Schulz, S.~P\'all, J.~C. Smith, B.~Hess, and
E.~Lindahl,
``GROMACS: high performance molecular simulations through multi-level
parallelism from laptops to supercomputers,''
\emph{SoftwareX} \textbf{1--2}, 19--25 (2015).
DOI: 10.1016/j.softx.2015.06.001.

\bibitem{shirts2017engines}
M.~R. Shirts, C.~Klein, J.~M. Swails, J.~Yin, M.~K. Gilson, D.~L. Mobley,
D.~A. Case, and E.~D. Zhong,
``Lessons learned from comparing molecular dynamics engines on the SAMPL5
dataset,''
\emph{J. Comput.-Aided Mol. Des.} \textbf{31}(1), 147--161 (2017).
DOI: 10.1007/s10822-016-9977-1.

\bibitem{braun2018bestpractices}
E.~Braun, J.~Gilmer, H.~B. Mayes, D.~L. Mobley, J.~I. Monroe, S.~Prasad, and
D.~M. Zuckerman,
``Best practices for foundations in molecular simulations [Article v1.0],''
\emph{Living J. Comput. Mol. Sci.} \textbf{1}(1), 5957 (2018).
DOI: 10.33011/livecoms.1.1.5957.

\bibitem{mosaics-software}
\emph{MOSAICS}, methodologies for optimization and sampling in computational
studies, P.~Minary, Department of Computer Science, University of Oxford.
\url{https://www.cs.ox.ac.uk/mosaics/} (accessed 8 September 2026). No
software paper exists for the engine; the underlying method is described by
Minary and Levitt (2010) and Sim, Levitt and Minary (2012), above.

\bibitem{pdb1efs}
V.~Larue and E.~Hantz,
``Conformation of a DNA-RNA hybrid,''
Protein Data Bank entry 1EFS, deposited 10 February 2000, released 6 March
2000; solution NMR, ten models.
DOI: 10.2210/pdb1efs/pdb.

\bibitem{hantz2001hybrid}
E.~Hantz, V.~Larue, P.~Ladam, L.~Le~Moyec, C.~Gouyette, and T.~Huynh~Dinh,
``Solution conformation of an RNA--DNA hybrid duplex containing a pyrimidine
RNA strand and a purine DNA strand,''
\emph{Int. J. Biol. Macromol.} \textbf{28}(4), 273--284 (2001).
DOI: 10.1016/S0141-8130(01)00123-4.

\bibitem{pdb1bna}
H.~R. Drew, R.~M. Wing, T.~Takano, C.~Broka, S.~Tanaka, K.~Itakura, and
R.~E. Dickerson,
``Structure of a B-DNA dodecamer: conformation and dynamics,''
\emph{Proc. Natl. Acad. Sci. U.S.A.} \textbf{78}(4), 2179--2183 (1981).
DOI: 10.1073/pnas.78.4.2179. Protein Data Bank entry 1BNA.

\bibitem{pdb2koc}
S.~Nozinovic, B.~F\"urtig, H.~R.~A. Jonker, C.~Richter, and H.~Schwalbe,
``High-resolution NMR structure of an RNA model system: the 14-mer cUUCGg
tetraloop hairpin RNA,''
\emph{Nucleic Acids Res.} \textbf{38}(2), 683--694 (2010).
DOI: 10.1093/nar/gkp956. Protein Data Bank entry 2KOC.

\bibitem{pdb1ubq}
S.~Vijay-Kumar, C.~E. Bugg, and W.~J. Cook,
``Structure of ubiquitin refined at 1.8~\AA{} resolution,''
\emph{J. Mol. Biol.} \textbf{194}(3), 531--544 (1987).
DOI: 10.1016/0022-2836(87)90679-6. Protein Data Bank entry 1UBQ.

\bibitem{pdb5pti}
A.~Wlodawer, J.~Walter, R.~Huber, and L.~Sj\"olin,
``Structure of bovine pancreatic trypsin inhibitor: results of joint neutron
and X-ray refinement of crystal form II,''
\emph{J. Mol. Biol.} \textbf{180}(2), 301--329 (1984).
DOI: 10.1016/S0022-2836(84)80006-6. Protein Data Bank entry 5PTI.

\bibitem{pdb1hhk}
D.~R. Madden, D.~N. Garboczi, and D.~C. Wiley,
``The antigenic identity of peptide-MHC complexes: a comparison of the
conformations of five viral peptides presented by HLA-A2,''
\emph{Cell} \textbf{75}(4), 693--708 (1993).
DOI: 10.1016/0092-8674(93)90490-H. Protein Data Bank entry 1HHK.

\bibitem{pdb1f7y}
E.~Ennifar, A.~Nikulin, S.~Tishchenko, A.~Serganov, N.~Nevskaya, M.~Garber,
B.~Ehresmann, C.~Ehresmann, S.~Nikonov, and P.~Dumas,
``The crystal structure of UUCG tetraloop,''
\emph{J. Mol. Biol.} \textbf{304}(1), 35--42 (2000).
DOI: 10.1006/jmbi.2000.4204. Protein Data Bank entry 1F7Y.

\bibitem{codata2018}
E.~Tiesinga, P.~J. Mohr, D.~B. Newell, and B.~N. Taylor,
``CODATA recommended values of the fundamental physical constants: 2018,''
\emph{Rev. Mod. Phys.} \textbf{93}(2), 025010 (2021).
DOI: 10.1103/RevModPhys.93.025010.

\bibitem{dolinsky2004pdb2pqr}
T.~J. Dolinsky, J.~E. Nielsen, J.~A. McCammon, and N.~A. Baker,
``PDB2PQR: an automated pipeline for the setup of Poisson--Boltzmann
electrostatics calculations,''
\emph{Nucleic Acids Res.} \textbf{32}(Web Server), W665--W667 (2004).
DOI: 10.1093/nar/gkh381.

\bibitem{jurrus2018apbs}
E.~Jurrus, D.~Engel, K.~Star, K.~Monson, J.~Brandi, L.~E. Felberg,
D.~H. Brookes, L.~Wilson, J.~Chen, K.~Liles, M.~Chun, P.~Li, D.~W. Gohara,
T.~Dolinsky, R.~Konecny, D.~R. Koes, J.~E. Nielsen, T.~Head-Gordon, W.~Geng,
R.~Krasny, G.-W. Wei, M.~J. Holst, J.~A. McCammon, and N.~A. Baker,
``Improvements to the APBS biomolecular solvation software suite,''
\emph{Protein Sci.} \textbf{27}(1), 112--128 (2018).
DOI: 10.1002/pro.3280.

\bibitem{fey2019pyg}
M.~Fey and J.~E. Lenssen,
``Fast graph representation learning with PyTorch Geometric,''
arXiv:1903.02428 (2019). \url{https://arxiv.org/abs/1903.02428}.

\bibitem{wang2019dgl}
M.~Wang, D.~Zheng, Z.~Ye, Q.~Gan, M.~Li, X.~Song, J.~Zhou, C.~Ma, L.~Yu,
Y.~Gai, T.~Xiao, T.~He, G.~Karypis, J.~Li, and Z.~Zhang,
``Deep Graph Library: a graph-centric, highly-performant package for graph
neural networks,''
arXiv:1909.01315 (2019). \url{https://arxiv.org/abs/1909.01315}.

\bibitem{rego2015threedmol}
N.~Rego and D.~Koes,
``3Dmol.js: molecular visualization with WebGL,''
\emph{Bioinformatics} \textbf{31}(8), 1322--1324 (2015).
DOI: 10.1093/bioinformatics/btu829.

\bibitem{zakai2011emscripten}
A.~Zakai,
``Emscripten: an LLVM-to-JavaScript compiler,''
in \emph{Proceedings of the ACM International Conference Companion on Object
Oriented Programming Systems Languages and Applications Companion (OOPSLA
'11)}, 301--312 (2011).
DOI: 10.1145/2048147.2048224. Version 6.0.9 was used.

\bibitem{pyodide-software}
\emph{Pyodide}, version 0.27.7, the Pyodide contributors.
\url{https://pyodide.org/} (accessed 8 September 2026). Cited as software; no
paper describes it.

\bibitem{hatch2025mcbestpractices}
H.~W. Hatch, D.~S. Corti, D.~A. Kofke, and V.~K. Shen,
``Best practices for developing Monte Carlo methodologies in molecular
simulations [Article v1.0],''
\emph{Living J. Comput. Mol. Sci.} \textbf{6}(1), 3289 (2025).
DOI: 10.33011/livecoms.6.1.3289.

\bibitem{lindorfflarsen2012validation}
K.~Lindorff-Larsen, P.~Maragakis, S.~Piana, M.~P. Eastwood, R.~O. Dror, and
D.~E. Shaw,
``Systematic validation of protein force fields against experimental data,''
\emph{PLoS ONE} \textbf{7}(2), e32131 (2012).
DOI: 10.1371/journal.pone.0032131.

\bibitem{beauchamp2012benchmark}
K.~A. Beauchamp, Y.-S. Lin, R.~Das, and V.~S. Pande,
``Are protein force fields getting better? A systematic benchmark on 524
diverse NMR measurements,''
\emph{J. Chem. Theory Comput.} \textbf{8}(4), 1409--1414 (2012).
DOI: 10.1021/ct2007814.

\bibitem{best2012charmm36}
R.~B. Best, X.~Zhu, J.~Shim, P.~E.~M. Lopes, J.~Mittal, M.~Feig, and
A.~D. MacKerell,~Jr.,
``Optimization of the additive CHARMM all-atom protein force field targeting
improved sampling of the backbone $\phi$, $\psi$ and side-chain $\chi_1$ and
$\chi_2$ dihedral angles,''
\emph{J. Chem. Theory Comput.} \textbf{8}(9), 3257--3273 (2012).
DOI: 10.1021/ct300400x.

\bibitem{huang2017charmm36m}
J.~Huang, S.~Rauscher, G.~Nawrocki, T.~Ran, M.~Feig, B.~L. de~Groot,
H.~Grubm\"uller, and A.~D. MacKerell,~Jr.,
``CHARMM36m: an improved force field for folded and intrinsically disordered
proteins,''
\emph{Nat. Methods} \textbf{14}(1), 71--73 (2017).
DOI: 10.1038/nmeth.4067.

\end{thebibliography}
